% !TeX root = ../main.tex

% ==================================================
% D.2. Series·DataFrame·Index의 객체 모델
% ==================================================

\subsection{Series·DataFrame·Index의 객체 모델}

\subsubsection{이전 학습과의 연결 및 핵심 질문}

\zsh{D.1}에서는 NumPy의 \zsh{ndarray}가 차원과 위치를 중심으로 데이터를 관리하고, Pandas의 \zsh{Series}와 \zsh{DataFrame}은 레이블을 포함한 데이터 구조를 제공한다는 차이를 배웠다.

이번에는 그 차이가 실제 Python 객체 모델에서 어떻게 구현되는지 살펴볼 것이다.

다음 코드를 생각해 보자.

\begin{dhkcode}
import pandas as pd

df = pd.DataFrame({
    "name": ["Kim", "Lee"],
    "score": [90, 85],
})

s = df["score"]
\end{dhkcode}

우리가 이미 알고 있는 Python 객체 모델에 따르면, \zsh{df}는 DataFrame 인스턴스를 참조하는 이름이고 \zsh{df['score']}는 \zsh{__getitem__()}에 의한 연산이다.

그런데 여기서 몇 가지 질문이 발생한다.

\begin{dhksmry}
1. DataFrame은 내부적으로 Series 객체들을 저장하고 있는 것인가?

2. df.index와 df.columns는 단순한 리스트일까, 별도의 객체일까?

3. df['score']로 반환된 Series는 원본 데이터와 어떤 관게를 가질까?

4. shape, ndim, dtype, dtypes는 어디에 저장되고 어떻게 조회될까?

5. DataFrame의 속성에 접근하는 것과 열을 선택하는 것은 동일한 연산인가?
\end{dhksmry}

이번 \zsh{D.2}의 목표는 Pandas 객체의 논리적 구조와 Python 객체로서의 실제 동작을 구분하는 것이다.

\subsubsection{Pandas의 세 가지 핵심 객체}

Pandas에서 이번 학습의 중심이 되는 클래스는 \zsh{Series}, \zsh{DataFrame}, \zsh{Index}다.

\begin{dhktable}
|클래스|객체의 역할|차원|
|:---|:---|:---|
|pd.Series|값과 행 레이블을 관리하는 1차원 데이터 구조|1|
|pd.DataFrame|행·열 레이블과 열별 dtype을 관리하는 2차원 데이터 구조|2|
|pd.Index|축에 대응하는 레이블의 순서와 조회·정렬 동작을 관리하는 객체|일반적으로 1차원|
\end{dhktable}

여기서 \zsh{Index}는 다른 두 객체와 역할이 조금 다르다.

Series와 DataFrame은 실제 관측값을 다루는 데이터 컨테이너다.

반면 Index는 어떤 값이 어느 축의 레이블에 대응하는지 관리하는 구조다.

\begin{dhknote}
\begin{dhksmry}[centered]
독립 실행 코드 1 - 클래스와 인스턴스
\end{dhksmry}
\end{dhknote}

\begin{dhkcode}
import pandas as pd

idx = pd.Index(["doc-A", "doc-B"])

s = pd.Series(
    [90, 85],
    index=idx,
    name="score",
)

df = pd.DataFrame({
    "score": s,
})

print(type(pd.Series))
print(type(s).__name__)
print(type(df).__name__)
print(type(idx).__name__)

print(isinstance(s, pd.Series))
print(isinstance(df, pd.DataFrame))
print(isinstance(idx, pd.Index))
\end{dhkcode}

예상 출력:

\begin{dhkzsh}
<class 'type'>
Series
DataFrame
Index
True
True
True
\end{dhkzsh}

\zsh{pd.Series}는 클래스 객체다.

\zsh{pd.Series(...)}는 클래스 객체를 호출하여 Series 인스턴스를 생성하는 표현식이다.

동일하게 \zsh{pd.DataFrame(...)}과 \zsh{pd.Index(...)}도 각각 해당 클래스의 인스턴스를 생성한다.

이 구조는 \zsh{class}, \zsh{type}, 인스턴스 생성 및 이름 바인딩 원칙과 동일하다.

하지만 중요한 차이가 있다.

Pandas가 제공하는 데이터 처리 규칙은 이 객체들의 메서드와 연산 프로토콜에 의해 결정된다는 점이다.

\subsubsection{Series의 객체 구조 - 값, Index, dtype, name}

Series는 1차원 데이터 구조지만 단순한 값의 나열만 관리하는 것은 아니다.

논리적으로 다음 네 요소를 구분할 수 있다.

\begin{dhkgraph}[uniform]
node n1 (1,3) [accent] :: Series 인스턴스 // s
node n2 (2,1) :: Values // 90, 85
node n3 (2,2) :: Index // doc-A, doc-B
node n4 (2,4) :: dtype // int 64
node n5 (2,5) :: name // score

split n1 -> n2,n3,n4,n5

explain :: 하나의 Series 객체가 제공하는 논리적 구성 // 실제 메모리 레이아웃과 동일한 그림은 아님
\end{dhkgraph}

이 그림에서 \zsh{Values}, \zsh{Index}, \zsh{dtype}, \zsh{name}을 모두 같은 종류의 독립적인 저장 객체라고 생각하면 안 된다.

값의 저장 표현은 NumPy 배열 또는 Pandas ExtensionArray 등일 수 있고, \zsh{dtype}은 그 값들의 자료형을 기술하는 객체다.

\zsh{name}은 Series에 부여한 이름으로, 데이터 값과는 별개다.

\begin{dhknote}
\begin{dhksmry}[centered]
독립 실행 코드 2 - 속성과 객체 동일성
\end{dhksmry}
\end{dhknote}

\begin{dhkcode}
import pandas as pd

idx = pd.Index(["doc-A", "doc-B"])

s = pd.Series(
    [90, 85],
    index=idx,
    name="score",
    dtype="int64",
)

print(s.shape)
print(s.ndim)
print(s.size)
print(s.dtype)
print(s.name)

print(s.index.tolist())
print(s.index is idx)
\end{dhkcode}

예상 출력:

\begin{dhkzsh}
(2,)
1
2
int64
score
['doc-A', 'doc-B']
True
\end{dhkzsh}

마지막마 \zsh{True}는 이 예제에서 생성자에게 전달한 \zsh{idx} 객체가 Series의 Index로 재사용되었음을 나타낸다.

그러나 이것을 ㅁ도느 생성 경로에서 Pandas가 반드시 동일한 Index 객체를 재사용한다는 일반 규칙으로 확대하면 안 된다.

생성자 옵션과 데이터 변환에 따라 다른 Index 객체가 만들어질 수도 있다.

\begin{dhknote}
\begin{dhksmry}[centered]
shape와 size의 차이
\end{dhksmry}
\end{dhknote}

\zsh{shape}은 튜플이고, \zsh{size}는 원소 개수를 나타내는 정수다.

\begin{dhkcode}
print(type(s.shape))
print(type(s.size))
\end{dhkcode}

예상 출력:

\begin{dhkzsh}
<class 'tuple'>
<class 'int'>
\end{dhkzsh}

따라서 \zsh{shape}은 배열의 구조를, \zsh{size}는 원소의 수를 나타내는 서로 다른 개념이다.

또한 \zsh{s.dtype}은 Series의 값에 대한 \zsh{dtype}을 반환한다.

\zsh{s.index.dtype}은 Index 레이블의 dtype이므로 서로 같을 필요가 없다.

\subsubsection{DataFrame의 객체 구조 - 두 개의 축과 열별 dtype}

Series가 하나의 Index를 갖는다면, DataFrame은 두 개의 축을 가진다.

\begin{dhksmry}
    - index : 행 축의 레이블

    - columns : 열 축의 레이블
\end{dhksmry}

그리고 DataFrame은 각 열마다 서로 다른 dtype을 가질 수 있다.

다음 예제를 살펴보자.

\begin{dhkcode}
import pandas as pd

df = pd.DataFrame(
    {
        "text": ["hello", "world"],
        "label": [1, 0],
        "score": [0.9, 0.7],
    },
    index=["doc-A", "doc-B"],
).astype({
    "text": "object",
    "label": "int64",
    "score": "float64",
})

print(df)
print(df.shape)
print(df.ndim)
print(df.dtypes)
\end{dhkcode}

예상 출력:

\begin{dhkzsh}
        text  label  score
doc-A  hello      1    0.9
doc-B  world      0    0.7

(2, 3)
2
text      object
label      int64
score    float64
dtype: object
\end{dhkzsh}

여기서 \zsh{df.shape == (2, 3)}이라는 사실은 DataFrame에 두 개의 행과 세 개의 열이 있다는 뜻이다.

그런데 \zsh{df.dtypes}는 왜 \zsh{Series}처럼 출력될까?

실제로 반환 객체가 Series이기 때문이다.

\begin{dhkcode}
print(type(df.dtypes).__name__)
print(df.dtypes.index.tolist())
print(df.dtypes.name)
\end{dhkcode}

예상 출력:

\begin{dhkzsh}
Series
['text', 'label', 'score']
None
\end{dhkzsh}

이것은 중요한 객체 관계를 보여준다.

\zsh{df.dtypes}는 DataFrame의 각 열에 대응하는 dtype을 하나의 Series로 표현한다.

즉, dtype 정보를 담은 Series의 Index가 원본 DataFrame의 열 레이블과 대응하는 것이다.

\begin{dhknote}
\begin{dhksmry}[centered]
DataFrame은 Series 객체를 저장하는 딕셔너리일까?
\end{dhksmry}
\end{dhknote}

다음과 같은 개념적 표현은 유용하다.

\begin{dhkdir}[in=6]
DataFrame
> "text" -> Series
> "label" -> Series
> "score" -> Series
\end{dhkdir}

하지만 이것을 실제 저장 구조라고 단정하면 잘못된 이해가 된다.

DataFrame의 내부 데이터 관리 시스템은 여러 열의 데이터를 블록 단위로 관리하거나 확장 배열(ExtensionArray)을 활용할 수 있다.

따라서 \zsh{df["label"]}이 Series를 반환한다는 사실만으로 DataFrame 내부에 해당 Series 인스턴스가 독립적으로 저장되어 있다고 결론 내릴 수 없다.

이전 \zsh{nn.Module}과 비교해 보면 차이가 분명해진다.

PyTorch에서는 \zsh{model._modules["linear"]}에 실제 하위 모듈 인스턴스가 등록되어 있었다.

반면 Pandas의 DataFrame을 구성하는 열은 반드시 \zsh{_modules}와 같은 딕셔너리에 Series 인스턴스로 등록되는 구조가 아니다.

열을 Series로 조회할 수 있다는 것과 열을 Series 인스턴스로 저장한다는 것은 다른 명제다.

\subsubsection{Index}는 어떤 객체이며, 왜 중요한가?

Index는 단순히 출력 화면 왼쪽에 표시되는 행 번호가 아니다.

Pandas는 Index를 통해 축의 레이블을 관리하고, 레이블 기반 조회와 정렬 연산을 수행할 수 있다.

Index 자체도 Python 객체이기 때문에 타입, dtype, 이름, 유일성 등을 확인할 수 있다.

\begin{dhknote}
\begin{dhksmry}[centered]
독립 실행 코드 3 - Index 객체 분석
\end{dhksmry}
\end{dhknote}

\begin{dhkcode}
import pandas as pd

idx = pd.Index(
    ["doc-A", "doc-B", "doc-C"],
    name="doc_id",
)

print(type(idx).__name__)
print(idx.ndim)
print(idx.shape)
print(idx.name)
print(idx.is_unique)
print(idx.has_duplicates)
print(idx.get_loc("doc-B"))
\end{dhkcode}

예상 출력:

\begin{dhkzsh}
Index
1
(3,)
doc_id
True
False
1
\end{dhkzsh}

여기서 \zsh{idx.get_loc("doc-B")}는 레이블을 기준으로 위치 정보를 찾는 메서드다.

현재 Index에서는 \zsh{"doc-B"}가 한 번만 등장하므로 정수 위치 1을 반환한다.

그런데 Index에 중복 레이블이 있다면 겨로가는 달라질 수 있다.

\begin{dhkcode}
import pandas as pd

idx = pd.Index(
    ["doc-A", "doc-B", "doc-B", "doc-C"]
)

print(idx.is_unique)
print(idx.get_loc("doc-B"))
\end{dhkcode}

예상 출력:

\begin{dhkzsh}
False
slice(1, 3, None)
\end{dhkzsh}

중복된 레이블이 연속해서 등장하기 때문에 해당 위치 범위를 표현하는 \zsh{slice}가 반환된 것이다.

여기서 중요한 점은 \zsh{get_loc()}이 언제나 정수를 반환하는 메서드가 아니라는 사실이다.

Index의 유일성고 중복 레이블의 배치에 따라 정수, 슬라이스 또는 불리언 마스크를 반환할 수 있다.

\begin{dhknote}
\begin{dhksmry}[centered]
Index의 불변성과 DataFrame의 변경 가능성
\end{dhksmry}
\end{dhknote}

Index는 일반적인 리스트처럼 원소를 직접 변경하는 방식을 지원하지 않는다.

\begin{dhkcode}
import pandas as pd

idx = pd.Index(["A", "B"])

try:
    idx[0] = "C"
except TypeError as e:
    print(type(e).__name__)
\end{dhkcode}

예상 출력:

\begin{dhkzsh}
TypeError
\end{dhkzsh}

그렇다고 DataFrame의 Index를 교체할 수 없다는 뜻은 아니다.

\begin{dhkcode}
import pandas as pd

df = pd.DataFrame(
    {"score": [90, 85]},
    index=["A", "B"],
)

old_index = df.index

df.index = ["C", "D"]

print(df.index.tolist())
print(old_index.tolist())
print(df.index is old_index)
\end{dhkcode}

예상 출력:

\begin{dhkzsh}
['C', 'D']
['A', 'B']
False
\end{dhkzsh}

여기에서는 기존 Index 객체의 원소를 수정한 것이 아니다.

DataFrame이 사용하는 행 레이블을 새 Index 구조로 교체한 것이다.

이 차이를 객체 모델로 표현하면 다음과 같다.

\begin{dhkmatrix}[delim=none,style=cards]
cols :: 변경 전, 변경 후
data :: [[df,df],[↓, ↓],[Index 객체 1 - "A" "B",green:Index 객체 2 - "C" "D"]]
explain :: 기존 Index 객체의 원소를 변경한 것이 아니라 DataFrame이 참조하는 축 레이블 객체를 교체한 상황
\end{dhkmatrix}

Index의 불변성은 기본적인 원소 변경 인터페이스에 관한 것이다.

내부에 변경 가능한 Python 객체가 들어 있는 특수한 경우까지 깊은 불변성(deep immutability)이 보장된다고 해석해서는 안 된다.

\subsubsection{DataFrame}의 열 선택은 어떤 객체를 반환하는가?

이제 Python의 특수 메서드와 연결해 보자.

\begin{dhkcode}
s = df["score"]
\end{dhkcode}

이 표현식은 DataFrame 객체에 문자열 키를 전달하는 구독(subscription) 연산이다.

일반적인 Python 규칙에서 \zsh{obj[key]}는 

객체 타입이 제공하는 \zsh{__getitem__()} 메커니즘으로 처리된다.

Pandas는 여기에 열 레이블 조회라는 의미를 부여한다.

\begin{dhknote}
\begin{dhksmry}[centered]
독립 실행 코드 4 - 반환 객체와 레이블
\end{dhksmry}
\end{dhknote}

\begin{dhkcode}
import pandas as pd

df = pd.DataFrame(
    {
        "score": [90, 85],
        "label": [1, 0],
    },
    index=["doc-A", "doc-B"],
)

a = df["score"]
b = df[["score"]]

print(type(a).__name__)
print(a.shape)
print(a.ndim)

print(type(b).__name__)
print(b.shape)
print(b.ndim)

print(a.index.tolist())
print(b.columns.tolist())
\end{dhkcode}

예상 출력:

\begin{dhkzsh}
Series
(2,)
1
DataFrame
(2, 1)
2
['doc-A', 'doc-B']
['score']
\end{dhkzsh}

이 결과는 매우 중요하다.

\zsh{df["score"]}와 \zsh{df[["score"]]}는 비슷한 데이터를 선택하지만 반환 객체의 타입이 다르다.

\begin{dhktable}
|표현식	|반환 타입|	shape|	ndim|
|:---|:---|:---|:---|
|df["score"]|	Series	|(2,)	|1|
|df[["score"]]|	DataFrame|	(2, 1)|	2|
\end{dhktable}

왜 이런 차이가 생길까?

첫 번째 표현식은 문자열 하나를 열 선택 키로 전달한다.

일반적인 유일한 열 레이블을 가진 DataFrame에서는 Series가 반환된다.

두 번째 표현식은 문자열 하나를 원소로 갖는 리스트를 전달한다.

따라서 여러 열을 선택하는 인터페이스를 사용하고, 결과도 2차원 DataFrame으로 유지된다.

\begin{dhknote}
\begin{dhksmry}[centered]
경계 조건 - 중복 열 이름
\end{dhksmry}
\end{dhknote}

여기에는 예외가 있다.

\begin{dhkcode}
import pandas as pd

df = pd.DataFrame(
    [[90, 85]],
    columns=["score", "score"],
)

result = df["score"]

print(type(result).__name__)
print(result.shape)
\end{dhkcode}

동일한 열 레이블이 두 번 존재하기 때문에 \zsh{df["score"]}는 두 열을 포함하는 DataFrame을 반환한다.

따라서 \zsh{df["열이름"]}은 무조건 Series를 반환한다고 암기하면 안 된다.

반환 타입을 판단할 때는 선택 키의 형태뿐 아니라 열 레이블의 유일성도 고려해야 한다.

이부분은 \zsh{D.4}에서 \zsh{.loc}, \zsh{.iloc}과 비교하며 더 깊게 분석할 것이다.

\subsubsection{속성 접근과 열 선택은 같은 연산인가?}

다음 두 표현식은 같은 값을 반환할 수 있다.

\begin{dhkcode}
df.score
df["score"]
\end{dhkcode}

하지만 내부적으로 동일한 문법은 아니다.

첫 번재는 Python의 속성 접근(attribute access)이다.

두 번째는 구독 연산(subscription)이다.

\begin{dhknote}
\begin{dhksmry}[centered]
독립 실행 코드 5 - 이름 충돌
\end{dhksmry}
\end{dhknote}

\begin{dhkcode}
import pandas as pd

df = pd.DataFrame({
    "score": [90, 85],
    "shape": [10, 20],
})

print(type(df.score).__name__)
print(df["score"].tolist())

print(df.shape)
print(df["shape"].tolist())
\end{dhkcode}

예상 출력:

\begin{dhkzsh}
Series
[90, 85]
(2, 2)
[10, 20]
\end{dhkzsh}

왜 \zsh{df.shape}은 \zsh{"shape"} 열을 반환하지 않을까?

\zsh{shape}은 이미 DataFrame이 제공하는 속성 이름이기 때문이다.

일반적인 속성 탐색에서 shape이라는 속성을 찾을 수 있으므로 Pandas의 열 이름 접근 편의 기능이 적용되지 않는다.

반면 \zsh{df["shape"]}은 열 레이블을 명시적으로 전달하는 표현식이기 때문에 해당 열을 반환한다.

이것은 속성 탐색 원칙과 연결된다.

Pandas의 점 표기법은 일반적인 Python 속성 탐색을 기반으로 하며, 일부 조건에서 열 레이블 접근을 편의 기능으로 지원한다.

모든 열 이름을 속성처럼 사용할 수 있는 것은 아니다.

따라서 분석 파이프라인에서는 열 선택에 \zsh{df["column"]} 형태를 기본으로 사용하는 것이 안전하다.

특히 열 이름이 메서드나 속성 이름과 충돌하거나 Python 식별자로 적합하지 않은 경우에 차이가 분명해진다.

\subsubsection{Series와 DataFrame의 객체 동일성 및 데이터 공유}

지금까지는 \zsh{df["score"]}가 Series를 반환한다는 사실을 확인했다.

그렇다면 반환된 Series는 DataFrame의 데이터를 복사한 새로운 객체일까, 아니면 기존 데이터를 공유하는 객체일까?

이 질문에는 두 가지 개념이 포함되어 있다.

\begin{dhksmry}
    1. Python 객체의 동일성(Identity)

    2. 내부 데이터 버퍼의 공유(memory sharing)
\end{dhksmry}

두 개념을 서로 다르다.

\begin{dhknote}
\begin{dhksmry}[centered]
독립 실행 코드 6 - 객체 동일성과 원본 변경
\end{dhksmry}
\end{dhknote}

\begin{dhkcode}
import pandas as pd

df = pd.DataFrame({
    "score": [90, 85],
})

s = df["score"]

print(s is df)
print(s.iloc[0])
print(df.iloc[0, 0])

s.iloc[0] = 100

print(s.iloc[0])
print(df.iloc[0, 0])
\end{dhkcode}

Pandas 3.x의 Copy-on-Write 동작을 기준으로 한 예상 출력은 다음과 같다.

\begin{dhkzsh}
False
90
90
100
90
\end{dhkzsh}

이 결과를 순서대로 분석해 보자.

첫째, \zsh{s is df}는 \zsh{False}다.

Series와 DataFrame은 서로 다른 Python 객체이기 때문이다.

둘째, Series를 선택하는 과정에서 내부 데이터를 공유할 수 있다.

\zsh{df["score"]}가 새로운 Series 객체를 반환한다고 해서 그 객체의 모든 데이터가 즉시 복사된다는 뜻은 아니다.

셋째, Pandas 3.x의 Copy-on-Write는 공유 데이터의 변경을 격리한다.

\zsh{s.iloc[0] = 100}을 수행할 때 필요한 경우 데이터 복사가 발생하여 원본 DataFrame의 값이 변경되지 않도록 한다.

Pandas 2.x에서는 설정과 선택 경로에 따라 동작이 달라질 수 있다.

따라서 이 예제의 변경 결과를 모든 버전에 적용해서는 안 된다.

여기서 이전 Python 객체 모델과 연결되는 중요한 사실이 있다.

\zsh{is} 연산자로 확인할 수 있는 것은 Python 객체의 동일성이지, 객체 내부의 메모리 버퍼 공유 여부가 아니다.

서로 다른 Series 객체라도 같은 데이터를 공유할 수 있고, 동일한 Series 객체라도 변경 연산 이후 내부 데이터 저장 공간이 교체될 수 있다.

이 문제는 \zsh{D.7}에서 다음 개념들과 함께 정식으로 분석할 것이다.

\begin{dhksmry}
    - 얕은 복사(shallow copy)와 깊은 복사(deep copy)

    - Copy-on-Write

    - 메서드 반환 객체와 원본 변경

    - 체인 인덱싱(chained indexing)

    - inplace 인수의 의미와 한계
\end{dhksmry}

\subsubsection{객체의 속성과 메서드를 구분하는 방법}

함수 객체와 함수 호출 결과를 구분했던 것을 기억해 보자.

Pandas에서도 이 구분이 중요하다.

\begin{dhkcode}
df.head
df.head()
\end{dhkcode}

첫 번재는 메서드 객체를 조회하는 표현식이고, 두 번째는 메서드를 호출하는 표현식이다.

반면 다음은 조금 다르다.

\begin{dhkcode}
df.shape
df.dtypes
\end{dhkcode}

이들은 메서드가 아니라 속성으로 제공된다.

\begin{dhknote}
\begin{dhksmry}[centered]
독립 실행 코드 7 - 속성, 메서드, 바인딩
\end{dhksmry}
\end{dhknote}

\begin{dhkcode}
import pandas as pd

df = pd.DataFrame({
    "score": [90, 85],
})

method = df.head
result = method(1)

print(callable(method))
print(method.__self__ is df)
print(method.__func__ is pd.DataFrame.head)

print(type(result).__name__)
print(result.shape)

print(callable(df.shape))
print(type(df.shape).__name__)
\end{dhkcode}

\begin{dhkzsh}
True
True
True
DataFrame
(1, 1)
False
tuple
\end{dhkzsh}

이 결과에서 확인할 수 있는 객체 관계는 다음과 같다.

\zsh{df.head}는 DataFrame 인스턴스에 바인딩된 메서드(bound method)다.

따라서 \zsh{method.__self__}는 \zsh{df}를 참조하고,

\zsh{method.__func__}는 해당 메서드의 기반 함수 객체를 참조한다.

반면 \zsh{df.shape}는 튜플을 반환하는 속성이므로 호출 가능한 메서드가 아니다.

\begin{dhknote}
\begin{dhksmry}[centered]
속성이라고 해서 단순히 __dict__에 저장된 값일까?
\end{dhksmry}
\end{dhknote}

그렇지는 않다.

Python의 속성은 클래스에 정의된 디스크립터(descriptor)를 통해 계산될 수도 있다.

예를 들어 \zsh{shape}는 DataFrame의 현재 축 길이를 바탕으로 제공되는 속성이다.

따라서 다음과 같은 등가 관계를 가정하면 안 된다.

\begin{dhkcode}
# 일반적으로 성립한다고 보장할 수 없는 표현
df.shape == df.__dict__["shape"]
\end{dhkcode}

여기에서 중요한 판단 기준은 속성 접근 문법, 실제 저장 위치, 속성의 계산 방식은 서로 구분해야 한다는 것이다.

이것은 \zsh{nn.Module}의 \zsh{_modules} 등록 구조를 분석할 때 사용했던 원칙과도 연결된다.

\subsubsection{ML/Transofrmer 응용 - 샘플 ID를 유지하는 분석 테이블}

이제 세 객체를 하나의 실제 분석 흐름으로 통합해 보자.

Transformer를 이용해 문서 분류 모델을 학습했다고 가정한다.

모델의 예측 결과를 분석하려면 최소한 다음 정보가 필요하다.

\begin{dhksmry}
    - 문서 식별자

    - 원본 텍스트

    - 정답 레이블

    - 예측 레이블

    - 토큰 길이
\end{dhksmry}

이 정보를 DataFrame으로 관리할 수 있다.

\begin{dhknote}
\begin{dhksmry}[centered]
독립 실행 코드 8 - 오류 분석 테이블
\end{dhksmry}
\end{dhknote}

\begin{dhkcode}
import pandas as pd

df = pd.DataFrame(
    {
        "text": [
            "good movie",
            "bad service",
            "excellent product",
        ],
        "label": [1, 0, 1],
        "prediction": [1, 1, 0],
        "token_length": [2, 2, 2],
    },
    index=pd.Index(
        ["doc-A", "doc-B", "doc-C"],
        name="doc_id",
    ),
)

df["correct"] = (
    df["label"] == df["prediction"]
)

errors = df.loc[
    ~df["correct"]
]

print(df["correct"].tolist())
print(errors.index.tolist())
print(errors.shape)
print(errors["text"].tolist())
\end{dhkcode}

예상 출력:

\begin{dhkzsh}
[True, False, False]
['doc-B', 'doc-C']
(2, 5)
['bad service', 'excellent product']
\end{dhkzsh}

객체 및 실행 과정 추적

\begin{dhkgraph}[uniform]
node n1 (1,1) :: df // DataFrame 3행 x 4열
node n2 (2,1) :: df['label'] == df['prediction'] // 비교 연산 -> bool Series 반환
node n3 (3,1) :: df['correct'] = ... // DataFrame에 열 추가 -> 3행 x 5열
node n4 (4,1) [accent] :: df.loc[~df['correct']] // 불리언 선택 -> 새 DataFrame 반환
node n5 (5,1) :: error.index // ["doc-B", "doc-C"]

edge n1->n2
edge n2 -> n3
edge n3 -> n4
\end{dhkgraph}

이 코드에서는 다음과 같은 동작이 순서대로 발생한다.

\begin{dhknote}
\begin{dhksmry}[centered]
1단계 - DataFrame 생성
\end{dhksmry}
\end{dhknote}

문서별 정보를 담은 DataFrame 인스턴스가 생성되고 \zsh{df}에 바인딩된다.

\begin{dhknote}
\begin{dhksmry}[centered]
2단계 - 열 선택과 비교
\end{dhksmry}
\end{dhknote}

\zsh{df["label"]}과 \zsh{df["prediction"]}은 각각 Series를 반환한다.

두 Series에 \zsh{==} 연산을 적용하면 각 행의 레이블이 대응되는 불리언 Series가 생성된다.

\begin{dhknote}
\begin{dhksmry}[centered]
3단계 - 새 열 할당
\end{dhksmry}
\end{dhknote}

\zsh{df["correct"] = ...}는 DataFrame에 새로운 열을 추가한다.

이때 \zsh{df}라는 이름이 다른 \zsh{DataFrame}에 재바인딩되는 것이 아니라, 기존 DataFrame 객체의 열 구성이 변경된다.

\begin{dhknote}
\begin{dhksmry}[centered]
4단계 - 오류 행 선택
\end{dhksmry}
\end{dhknote}

\zsh{~df["correct"]}는 불리언 값을 반전한 Series를 반환한다.

\zsh{.loc[...]}는 이 불리언 조건에 대응하는 행을 선택해 새로운 DataFrame을 반환한다.

\begin{dhknote}
\begin{dhksmry}[centered]
5단계 - 문서 ID 보존
\end{dhksmry}
\end{dhknote}

선택 결과인 \zsh{errors}는 원본의 행 Index 레이블을 유지한다.

따라서 어떤 문서에서 오류가 발생했는지 식별할 수 있다.

이것은 단순히 행번호 1,2를 저장하는 것보다 유용하다.

원본 데이터가 정렬되거나 필터링되어도 문서 ID를 기반으로 다른 데이터와 대응시킬 수 있기 때문이다.

단, Index가 중복되어 있거나 다른 테이블과의 식별자 대응이 깨진 경우에는 별도의 정합성 검사가 필요하다.

\subsubsection{비교 및 예외 - 이번 단원에서 반드시 구분할 개념}

\begin{dhktable}
|구분	|올바른 이해|
|:---|:---|
|Series와 ndarray|	둘 다 1차원 값을 표현할 수 있지만 Series에는 Index와 Pandas 연산 의미론이 있음|
|DataFrame과 2차원 ndarray|	DataFrame은 두 축의 레이블과 열별 dtype을 관리|
|Index와 Python list|	Index는 레이블 조회·정렬 기능을 제공하며 일반적인 원소 직접 변경을 허용하지 않음|
|df["x"]와 df[["x"]]|	일반적으로 Series와 DataFrame을 반환하지만 중복 열 레이블 등 예외가 있음|
|df.x와 df["x"]|	속성 접근과 열 선택은 서로 다른 연산|
|df.shape와 df.dtypes|	각각 튜플과 Series를 반환하는 속성|
|새로운 Series 객체와 데이터 복사|	객체가 새로 생성되었다고 내부 데이터가 반드시 복사되는 것은 아님|
|Index 레이블과 기본키|	Index는 중복을 허용할 수 있으므로 자동으로 기본키가 되지 않음|
\end{dhktable}

특히 다음 세 가지를 이번 학습의 핵심 판단 기준으로 삼으면 된다.

첫째, 속성 접근 결과가 어떤 타입의 객체인지 확인해야 한다.

Series, DataFrame, Index, tuple, dtype 객체는 모두 서로 다른 의미를 가진다.

둘째, 반환된 객체가 원본과 어떤 관계를 갖는지 구분해야 한다.

반환 타입, 객체 동일성, 데이터 공유, 원본 변경 여부는 각각 별도로 판단해야 한다.

셋째, 데이터의 위치와 레이블을 혼동해서는 안 된다.

\zsh{Index}는 데이터의 의미를 연결하는 축이지만, 유일성과 정합성이 자동으로 보장되는 식별자 체계는 아니다.

\subsubsection{학습 정리 및 다음 학습}

\zsh{D.2}에서는 데이터 객체가 어떤 구조로 존재하는지 살펴보았다.

하지만 아직 중요한 질문이 남아 있다.

정수와 문자열, 실수와 결측값이 하나의 열에 섞이면 Pandas는 어떤 dtype을 선택할까?

\zsh{None}, \zsh{np.nan}, \zsh{pd.NA}, \zsh{pd.NaT}는 모두 결측을 표현하지만 왜 서로 다른 객체이며, 비교 연산에서는 어떤 차이를 보일까?

\zsh{D.3}에서는 이 질문을 시작으로 NumPy dtype, Pandas ExtensionDtype, Nullable 자료형, 결측값의 전파 규칙과 연산 예외까지 연결할 것이다.

공식 참고자료는 다음과 같다.
\begin{dhksmry}
https://pandas.pydata.org/docs/user_guide/dsintro.html
https://pandas.pydata.org/docs/user_guide/indexing.html
https://pandas.pydata.org/docs/user_guide/migration.html
\end{dhksmry}